\documentclass[12pt,a4paper]{article}

% ====== 基础宏包（尽量少）======
\usepackage{ctex}        % 中文支持
\usepackage{geometry}    % 页面设置
\usepackage{amsmath,amssymb} % 基本数学公式
\usepackage{graphicx}    % 插图
\usepackage{setspace}    % 行距

% ====== 页面设置 ======
\geometry{left=2.5cm,right=2.5cm,top=2.5cm,bottom=2.5cm}
\onehalfspacing

\begin{document}

% =========================
% 标题部分
% =========================
\begin{center}
    {\LARGE \textbf{文章题目}} \\[1em]
    % ↑ 在此填写论文题目

    学号1\ 姓名1,\ 学号2\ 姓名2,\ 学号3\ 姓名3
    % ↑ 在此填写作者信息（学号 + 姓名）
\end{center}

\vspace{1em}

% =========================
% 摘要与关键词
% =========================
\noindent\textbf{摘要：}
% ↑ 在此撰写摘要：
% 简要说明研究背景、研究方法、主要结果和结论
本文的主要思想和结果。

\vspace{0.5em}

\noindent\textbf{关键词：}
特征值；非负矩阵分解
% ↑ 在此填写关键词，中文分号隔开

\vspace{1.5em}

% =========================
% 正文开始
% =========================

\section{研究背景}
% ↑ 介绍研究背景、研究动机和相关工作
% 说明：做什么？为什么做？别人怎么做？本文打算怎么做？

本节介绍研究背景与研究意义。

换行的效果
\section{第二节标题自拟}
% ↑ 可介绍理论基础、模型假设或算法思想
% 本节可以拆分为若干小节

\subsection{小节标题示例}
% ↑ 在此撰写具体内容，如数学模型或算法推导

这里写具体内容。

\section{第三节标题自拟}
% ↑ 描述本文提出的主要方法、理论或改进算法

这里描述你自己设计的理论或算法。

\section{第四节标题自拟}
% ↑ 实验与结果分析
% 可描述实验环境、参数设置、对比结果等

这里给出实验结果与分析。

\section{结论与展望}
% ↑ 总结全文工作，并给出未来可进一步研究的方向

本文总结了主要工作，并对未来研究进行了展望。

% =========================
% 参考文献
% =========================
\begin{thebibliography}{99}

\bibitem{Duff1986}
I.~S. Duff, A.~M. Erisman, J.~K. Reid,
\textit{Direct Methods for Sparse Matrices},
Oxford University Press, 1986.

% ↑ 在此继续添加参考文献

\end{thebibliography}

\end{document}